% TOP_PROBLEM: {"id": 2750, "title": "Kirby Problem 2.2", "category_id": 7, "category": {"id": 7, "name": "topology", "display_name": "Topology", "description": "Properties preserved under continuous deformations.", "slug": "topology", "order_index": 7, "created_at": "2026-07-31T15:26:25.671Z"}, "status": "open", "problem_number": "KP-2.2", "set_id": 11}
% FIRST_POSTED: 2026-10-01
% ENTRY_KIND: statement_only
% UnsolvedMath ID 2750; source catalogue code KP-2.2
% !TeX program = lualatex
% Definition and sources only; this file is not an LLM solution attempt.
\documentclass[11pt,a4paper]{article}
\usepackage[margin=25mm]{geometry}
\usepackage{fontspec}
\setmainfont{lmroman10-regular.otf}[BoldFont=lmroman10-bold.otf,
 ItalicFont=lmroman10-italic.otf,BoldItalicFont=lmroman10-bolditalic.otf]
\usepackage{amsmath,amssymb,mathtools}
\usepackage{unicode-math}
\setmathfont{latinmodern-math.otf}
\AtBeginDocument{\let\setminus\smallsetminus}
\usepackage{xurl,hyperref}
\hypersetup{colorlinks=true,urlcolor=blue}
\setcounter{secnumdepth}{0}
\setlength{\parindent}{0pt}
\setlength{\parskip}{5pt}
\setlength{\emergencystretch}{3em}
\begin{document}
\section*{Kirby Problem 2.2}\label{2750}
{\small UnsolvedMath ID 2750 · KP-2.2 · Topology · Kirby's Problems in Low-Dimensional Topology}
\subsection{Definitions and mathematical statement}
Let $S$ be a connected oriented surface of finite type, with punctures allowed,
and write $\operatorname{Mod}(S)$ for its orientation-preserving mapping class
group: homeomorphisms modulo isotopy. In this formulation there are no boundary
circles fixed pointwise; equivalently, boundary components are treated as punctures.
Put $\pi=\pi_1(S)$. A subgroup $K\leq\pi$ is characteristic if every automorphism
of $\pi$ preserves $K$. For a finite-index characteristic $K$, the action of mapping
classes on $\pi$ gives a homomorphism
\[
 \rho_K:\operatorname{Mod}(S)\longrightarrow\operatorname{Out}(\pi/K),
 \qquad \Gamma(K)=\ker\rho_K.
\]
Here $\operatorname{Out}(Q)$ is the automorphism group of $Q$ modulo inner
automorphisms; using it removes the choice of basepoint. Call a subgroup
\emph{congruence} if it contains some $\Gamma(K)$.

The congruence subgroup problem asks whether every finite-index subgroup
$H\leq\operatorname{Mod}(S)$ is congruence. Thus the requested assertion is that
for every such $H$ there exists a finite-index characteristic subgroup $K$ of
$\pi$ with $\Gamma(K)\subseteq H$. The finite quotients $\pi/K$ may be arbitrary;
restricting to homology modulo an integer would be a different, stronger requirement.
\subsection{Short English statement}
Does every finite-index subgroup of a surface mapping class group contain the kernel of its action on a suitable finite quotient of the surface group?
\subsection{Sources}
\begin{itemize}
\item[S1] ULAM.AI and the UnsolvedMath Contributors, supplied problems.json, record 2750 (KP-2.2), Kirby's Problems in Low-Dimensional Topology. Catalogue identity and source transcription; CC BY 4.0.
\url{https://huggingface.co/datasets/ulamai/UnsolvedMath}.
\item[S2] K3: A New Problem List in Low-Dimensional Topology, Problem 2.2. Expanded formulation of the supplied catalogue statement.
\url{https://math.berkeley.edu/sites/default/files/surv-295-ruberman-watermarked-author-pdf.pdf}.
\end{itemize}
\end{document}
